% Cote 119: offprints, reports, courses, CV and letters (typed and printed part): archive cote 119, site page 172 (PDF page 173).
% typescript. Draft compared: none.
% AI-transcribed by Claude Opus 5.5 (claude-opus-5-5, adaptive thinking, effort high), Message Batches API, batch r119_full (msgbatch_01DwZNm5Y7YsKBWyAN4sFs8w), 2026-11-13; not yet checked by eye.
\scanpage{119/172}
\mspage{- 11 -}

\uline{Théorème 4.4}.

Si $f(Z)$ est continue dans un ouvert $U$ du plan complexe et holo-%
morphe en tout point de $U$ sauf éventuellement en des points isolés la forme
$f(Z)\ dZ$ est fermée.

\uline{Démonstration}. D'après le théorème de Bolzano-Weier\unsure{strass} si $R$ est un
rectangle contenu dans $U$ il ne peut y avoir qu'un nombre fini de points de
$R$ où $f$ n'est pas holomorphe. En utilisant un découpage de $R$ par des paral-%
lèles aux axes passant par ces points (découpage analogue à celui rencontré
dans 3.6.) on est amené à démontrer que $\int_{\gamma} f(Z)\ dZ = 0$ dans le cas où $\gamma$
est un rectangle contenant au plus un sommet où $f$ n'est pas holomorphe. Soit
$(x_0,y_0)$ un tel sommet, $(x_0+a,\ y_0)$, $(x_0,y_0+b)$, $(x_0+a,\ y_0+b)$ les autres sommets
du rectangle. Supposons par exemple $a > 0$, $b > 0$. D'après 4.3. l'intégrale
de $f(Z)\ dZ$ le long de $\gamma_{\varepsilon}$ bord du rectangle $(x_0,y_0+\varepsilon)$ , $(x_0+a,y_0+\varepsilon)$ ,
$(x_0+a,y_0+b)$, $(x_0,\ y_0+b)$ est nulle et du fait de la continuité de $f$ cette
intégrale, tend vers $\int_{\gamma} f(Z)\ dZ$ quand $\varepsilon$ tend vers 0. Ce qui nous donne le
résultat attendu.

5 . \uline{Développement en série entière d'une fonction holomorphe}

\uline{Définition 5.1}.

Etant donnés un chemin différentiable fermé $\gamma$ dans le plan $C$ et
$Z_0$ un point de $C$ n'appartenant pas à l'image de $\gamma$ on appelle indice de $\gamma$
par rapport à $Z_0$ et on note $I(\gamma,Z_0)$ la valeur de l'intégrale
\[ \frac{1}{2\Pi i} \int_{\gamma} \frac{dZ}{Z-Z_0} \]

\uline{Proposition 5.2.}

$I(\gamma,Z)$ est un entier constant quend $Z$ vraie dans une composante
connexe du complémentaire de l'image de $\gamma$. Pour $Z$ dans la composante non
bornée de ce complémentaire on a $I(\gamma,Z) = 0$ .

\margin{172}
